\begin{abstract}
\noindent
El presente documento aborda el problema de enrutamiento de camiones y remolques (\textit{Truck and Trailer Routing Problem}, TTRP), una extensión NP-difícil del problema de enrutamiento de vehículos clásico que incorpora restricciones de accesibilidad y capacidad mediante una flota con remolques desacoplables. El objetivo principal de este informe es revisar el estado del arte de las metodologías de resolución desarrolladas en la literatura, analizando los principales enfoques exactos, heurísticos y metaheurísticos planteados desde su formulación inicial. Asimismo, se presenta la caracterización del problema, sus variantes operativas y la formulación de un modelo matemático de programación entera. Finalmente, se analizan las limitaciones de los métodos actuales y se discuten oportunidades de desarrollo para la implementación de algoritmos de optimización.
\end{abstract}
\vspace{1.5em}

\section*{Resumen}
\textit{
El presente documento aborda el problema de enrutamiento de camiones y remolques (\textit{Truck and Trailer Routing Problem}, TTRP), una extensión NP-difícil del problema de enrutamiento de vehículos clásico que incorpora restricciones de accesibilidad y capacidad mediante una flota con remolques desacoplables. El objetivo principal de este informe es revisar el estado del arte de las metodologías de resolución desarrolladas en la literatura, analizando los principales enfoques exactos, heurísticos y metaheurísticos planteados desde su formulación inicial. Asimismo, se presenta la caracterización del problema, sus variantes operativas y la formulación de un modelo matemático de programación entera. Finalmente, se analizan las limitaciones de los métodos actuales y se discuten oportunidades de desarrollo para la implementación de algoritmos de optimización.
}
\vspace{1em}